% GENERATED FILE. Edit canonical lessons, then run npm run generate:books.
% canonical-source-hash: fnv1a64:e63501b8f85d7b80
% canonical-lessons: ZH-C248-xueqi, ZH-C248-nianji, ZH-C248-xiaozhang, ZH-W248-xiao-school, ZH-C248-gaozhong, ZH-W248-gao-tall, ZH-C248-jiazhang

\chapter{A New Term}
\label{ch:zh-a-new-term}
\hlchaptermodality{248}

\begin{chapteropening}
\textbf{By the end of this chapter:} \emph{I can talk about the school term, my year group, the head teacher, senior high school and parents.}

Complete the last lesson of chapter 248: I can talk about the school term, my year group, the head teacher, senior high school and parents.
\end{chapteropening}

\section[\texorpdfstring{\zh{学期} (xuéqī)}{学期 (xuéqī)}]{\zh{学期} (xuéqī) — school term; semester}

\label{lesson:ZH-C248-xueqi}



\begin{quote}
Before the new one: say the Mandarin for souvenir, then the Mandarin for to recall.
\end{quote}

\subsection*{You'll want to know: \zh{学期}}
\textbf{\zh{学期}} — \emph{xuéqī} — \textquotedblleft{}school term; semester\textquotedblright{}.

\begin{quote}
\textbf{\zh{这个学期很长}} \emph{zhège xuéqī hěn cháng} — This term is a long one.
\end{quote}

\begin{cousinweb}
\textbf{\zh{学期}} (\emph{xuéqī}) is built from \zh{学} 'to study' and \zh{期} 'a set time', both inherited from Middle Chinese, \emph{haewk} and \emph{gi} (Guangyun, in Baxter's transcription), which glosses the second 'to meet' and 'a limit'. Middle Chinese \emph{gi} would regularly give a rising tone, qí; the level tone of qī is irregular.
\end{cousinweb}

\begin{grammarlens}[title={What you've built}]
One more word for saying what I do.
\end{grammarlens}

\subsection*{Your turn}
\begin{itemize}
\raggedright
  \item \emph{Say it:} \emph{xuéqī}
  \item \emph{Say it:} \emph{xuéqī}, once more
  \item \emph{Say it:} \emph{zhège xuéqī hěn cháng} — This term is a long one.
  \item \emph{Recall:} say the Mandarin for souvenir, then the Mandarin for to recall, then say \emph{xuéqī} again
\end{itemize}

\begin{tcolorbox}[breakable,colback=teal!4,colframe=teal!35!black,title={Before you move on}]
What does \zh{学期} mean? (\textquotedblleft{}School term; semester\textquotedblright{}.) Say it once more.
\end{tcolorbox}
\section[niánjí]{niánjí — year; grade (at school)}

\label{lesson:ZH-C248-nianji}



\begin{quote}
Before the new one: say the Mandarin for to recall, then the Mandarin for a school term.
\end{quote}

\subsection*{You'll want to know: niánjí}
\emph{niánjí} — \textquotedblleft{}year; grade (at school)\textquotedblright{}.

\begin{quote}
\emph{wǒmen niánjí yǒu yìbǎi ge xuésheng} — There are a hundred students in our year.
\end{quote}

Its characters come later; for now it is a word for the ear.

\begin{cousinweb}
\emph{niánjí} is built from \zh{年} 'a year' and \emph{jí} 'a rank', both inherited from Middle Chinese, \emph{nen} and \emph{kip} (Guangyun, in Baxter's transcription), which glosses the second 'rank' and quotes an older dictionary for it, 'the order of silk threads'.
\end{cousinweb}

\begin{grammarlens}[title={What you've built}]
One more word for saying what I do.
\end{grammarlens}

\subsection*{Your turn}
\begin{itemize}
\raggedright
  \item \emph{Say it:} \emph{niánjí}
  \item \emph{Say it:} \emph{niánjí}, once more
  \item \emph{Say it:} \emph{wǒmen niánjí yǒu yìbǎi ge xuésheng} — There are a hundred students in our year.
  \item \emph{Recall:} say the Mandarin for to recall, then the Mandarin for a school term, then say \emph{niánjí} again
\end{itemize}

\begin{tcolorbox}[breakable,colback=teal!4,colframe=teal!35!black,title={Before you move on}]
What does niánjí mean? (\textquotedblleft{}Year; grade (at school)\textquotedblright{}.) Say it once more.
\end{tcolorbox}
\section[\texorpdfstring{\zh{校长} (xiàozhǎng)}{校长 (xiàozhǎng)}]{\zh{校长} (xiàozhǎng) — head teacher; principal}

\label{lesson:ZH-C248-xiaozhang}



\begin{quote}
Before the new one: say the Mandarin for a school term, then the Mandarin for a year group at school.
\end{quote}

\subsection*{You'll want to know: \zh{校长}}
\textbf{\zh{校长}} — \emph{xiàozhǎng} — \textquotedblleft{}head teacher; principal\textquotedblright{}.

\begin{quote}
\textbf{\zh{他是我们学校的校长}} \emph{tā shì wǒmen xuéxiào de xiàozhǎng} — He's the head of our school.
\end{quote}

The new character, in pieces: \zh{校} — ten strokes in two pieces: on the left the tree piece, four strokes, as on the left of \zh{机}; on the right a dot and a bar, two short strokes under them, then a long stroke falling to the left crossed by a long stroke falling to the right. The next lesson puts it in your hand.

\begin{cousinweb}
\textbf{\zh{校长}} (\emph{xiàozhǎng}) is built from \zh{校} 'a school' and \zh{长} 'the head', both inherited from Middle Chinese, \emph{haewH} and \emph{trjangX} (Guangyun, in Baxter's transcription). The dictionary glosses the second 'big'; it lists the first in the title of an army officer, and the sense 'school' is not given there, so how it came is not settled.
\end{cousinweb}

\begin{grammarlens}[title={What you've built}]
One more word for saying what I do.
\end{grammarlens}

\subsection*{Your turn}
\begin{itemize}
\raggedright
  \item \emph{Say it:} \emph{xiàozhǎng}
  \item \emph{Say it:} \emph{xiàozhǎng}, once more
  \item \emph{Say it:} \emph{tā shì wǒmen xuéxiào de xiàozhǎng} — He's the head of our school.
  \item \emph{Recall:} say the Mandarin for a school term, then the Mandarin for a year group at school, then say \emph{xiàozhǎng} again
\end{itemize}

\begin{tcolorbox}[breakable,colback=teal!4,colframe=teal!35!black,title={Before you move on}]
What does \zh{校长} mean? (\textquotedblleft{}Head teacher; principal\textquotedblright{}.) Say it once more.
\end{tcolorbox}
\section[\texorpdfstring{\zh{校} (xiào)}{校 (xiào)}]{\zh{校} — the tree piece, then a lid and a cross}

\label{lesson:ZH-W248-xiao-school}



\begin{quote}
The word first, then the character inside it.

\begin{itemize}
\raggedright
  \item \emph{Recall:} say \emph{xiàozhǎng} — \textbf{R1}, one lesson back, before it decays
\end{itemize}

\begin{quote}
\zh{机} \zh{学}
\end{quote}

\emph{Read it:} each one aloud, with nothing said first — \emph{jī}, \emph{xué}
\end{quote}

\subsection*{Script — the tree piece, then a lid and a cross}
\hlblockfigure{\detokenize{figures/ZH-W248-xiao-school-filmstrip.pdf}}{How \zh{校} is written, stroke by stroke}

\begin{quote}
\zh{校}
\end{quote}

\textbf{Two pieces, left before right.} Count the pieces, not the strokes.

\textbf{The left piece, the tree piece, four strokes.} A short bar. A vertical down through it. A stroke falling to the left from the crossing. Then a short dot to the right: on the left of a character the last stroke is a dot, not a long sweep.

\textbf{The right piece, six strokes.} A dot on top and a bar under it. Two short strokes under the bar, the left one falling to the left, the right one a dot. Then a long stroke falling to the left and, across it, a long stroke falling to the right to the foot.

4 + 6 = ten.

\subsection*{Your turn}
\begin{itemize}
\raggedright
  \item \emph{Write it:} \zh{校}, all ten, with the model in front of you
  \item \emph{Cover:} write it again from memory, and count the pieces, not the strokes
\end{itemize}

\begin{tcolorbox}[breakable,colback=teal!4,colframe=teal!35!black,title={Before you move on}]
How many strokes in \zh{校}? (\textbf{Ten}.) What is the left piece? (\textbf{The tree piece, four strokes, as in \zh{机}}.) What are the last two strokes? (\textbf{A long fall to the left, crossed by a long fall to the right}.)

Source: \href{https://raw.githubusercontent.com/chanind/hanzi-writer-data/68d10a4b21150cae5e1ebbd223eed289cf32d90c/data/\%E6\%A0\%A1.json}{Hanzi Writer Data \zh{校}, pinned stroke paths and medians}.
\end{tcolorbox}
\section[\texorpdfstring{\zh{高中} (gāozhōng)}{高中 (gāozhōng)}]{\zh{高中} (gāozhōng) — senior high school}

\label{lesson:ZH-C248-gaozhong}



\begin{quote}
Before the new one: say the Mandarin for a year group at school, then the Mandarin for a head teacher.
\end{quote}

\subsection*{You'll want to know: \zh{高中}}
\textbf{\zh{高中}} — \emph{gāozhōng} — \textquotedblleft{}senior high school\textquotedblright{}.

\begin{quote}
\textbf{\zh{我们学校有高中}} \emph{wǒmen xuéxiào yǒu gāozhōng} — Our school has a senior high section.
\end{quote}

The new character, in pieces: \zh{高} — ten strokes: a dot and a long bar, then \zh{口}; under them a frame open at the foot, with a small \zh{口} inside it. The next lesson puts it in your hand.

\begin{cousinweb}
\textbf{\zh{高中}} (\emph{gāozhōng}) is built from \zh{高} 'high' and \zh{中} 'middle', both inherited from Middle Chinese, \emph{kaw} and \emph{trjuwng} (Guangyun, in Baxter's transcription), which glosses the first 'up' and 'lofty': the higher middle school, a short form of a longer name.
\end{cousinweb}

\begin{grammarlens}[title={What you've built}]
One more word for saying what I do.
\end{grammarlens}

\subsection*{Your turn}
\begin{itemize}
\raggedright
  \item \emph{Say it:} \emph{gāozhōng}
  \item \emph{Say it:} \emph{gāozhōng}, once more
  \item \emph{Say it:} \emph{wǒmen xuéxiào yǒu gāozhōng} — Our school has a senior high section.
  \item \emph{Recall:} say the Mandarin for a year group at school, then the Mandarin for a head teacher, then say \emph{gāozhōng} again
\end{itemize}

\begin{tcolorbox}[breakable,colback=teal!4,colframe=teal!35!black,title={Before you move on}]
What does \zh{高中} mean? (\textquotedblleft{}Senior high school\textquotedblright{}.) Say it once more.
\end{tcolorbox}
\section[\texorpdfstring{\zh{高} (gāo)}{高 (gāo)}]{\zh{高} — a lid, \zh{口}, then a frame around a small \zh{口}}

\label{lesson:ZH-W248-gao-tall}



\begin{quote}
The word first, then the character inside it.

\begin{itemize}
\raggedright
  \item \emph{Recall:} say \emph{gāozhōng} — \textbf{R1}, one lesson back, before it decays
\end{itemize}

\begin{quote}
\zh{口} \zh{中}
\end{quote}

\emph{Read it:} each one aloud, with nothing said first — \emph{kǒu}, \emph{zhōng}
\end{quote}

\subsection*{Script — a lid, \zh{口}, then a frame around a small \zh{口}}
\hlblockfigure{\detokenize{figures/ZH-W248-gao-tall-filmstrip.pdf}}{How \zh{高} is written, stroke by stroke}

\begin{quote}
\zh{高}
\end{quote}

\textbf{Three pieces, top to bottom.}

\textbf{The lid, two strokes.} A dot, then a long bar across.

\textbf{The middle piece, \zh{口}, three strokes.} A short vertical, a stroke that runs right and turns down, then the bar that closes the foot.

\textbf{The bottom piece, five strokes.} A vertical down the left. A bar across the top that turns down the right side and hooks to the left at the foot. Then a small \zh{口} inside the frame, three strokes as before.

2 + 3 + 5 = ten.

\subsection*{Your turn}
\begin{itemize}
\raggedright
  \item \emph{Write it:} \zh{高}, all ten, with the model in front of you
  \item \emph{Cover:} write it again from memory, and count the pieces, not the strokes
\end{itemize}

\begin{tcolorbox}[breakable,colback=teal!4,colframe=teal!35!black,title={Before you move on}]
How many strokes in \zh{高}? (\textbf{Ten}.) How many \zh{口} does it hold? (\textbf{Two: one in the middle, a small one inside the frame}.) Is the frame closed at the foot? (\textbf{No, it is open}.)

Source: \href{https://raw.githubusercontent.com/chanind/hanzi-writer-data/68d10a4b21150cae5e1ebbd223eed289cf32d90c/data/\%E9\%AB\%98.json}{Hanzi Writer Data \zh{高}, pinned stroke paths and medians}.
\end{tcolorbox}
\section[\texorpdfstring{\zh{家长} (jiāzhǎng)}{家长 (jiāzhǎng)}]{\zh{家长} (jiāzhǎng) — pupil's parent; parent of a schoolchild}

\label{lesson:ZH-C248-jiazhang}



\begin{quote}
Before the new one: say the Mandarin for a head teacher, then the Mandarin for senior high school.
\end{quote}

\subsection*{You'll want to know: \zh{家长}}
\textbf{\zh{家长}} — \emph{jiāzhǎng} — \textquotedblleft{}pupil's parent; parent of a schoolchild\textquotedblright{}.

\begin{quote}
\textbf{\zh{今天家长来学校}} \emph{jīntiān jiāzhǎng lái xuéxiào} — The parents are coming to school today.
\end{quote}

\begin{cousinweb}
\textbf{\zh{家长}} (\emph{jiāzhǎng}) is built from \zh{家} 'a home' and \zh{长} 'the head', both inherited from Middle Chinese, \emph{kae} and \emph{trjangX} (Guangyun, in Baxter's transcription): the head of a household, and at school a pupil's parent.
\end{cousinweb}

\begin{grammarlens}[title={What you've built}]
One more word for saying what I do.
\end{grammarlens}

\subsection*{Your turn}
\begin{itemize}
\raggedright
  \item \emph{Say it:} \emph{jiāzhǎng}
  \item \emph{Say it:} \emph{jiāzhǎng}, once more
  \item \emph{Say it:} \emph{jīntiān jiāzhǎng lái xuéxiào} — The parents are coming to school today.
  \item \emph{Recall:} say the Mandarin for a head teacher, then the Mandarin for senior high school, then say \emph{jiāzhǎng} again
\end{itemize}

\begin{tcolorbox}[breakable,colback=teal!4,colframe=teal!35!black,title={Before you move on}]
What does \zh{家长} mean? (\textquotedblleft{}Pupil's parent; parent of a schoolchild\textquotedblright{}.) Say it once more.
\end{tcolorbox}
